\part{BLOQUE III: METODOLOGÍA Y CASO DE ESTUDIO: IMPUTACIÓN DE DATOS}

\chapter{Imputación de Datos en Series Temporales}
\label{chap:metodologia}
\glsresetall

% --- SECCIÓN 1: INTRODUCCIÓN AL PROBLEMA ---
%\section{El Problema de los Datos Faltantes}
\label{sec:intro_datos_faltantes}
Los conjuntos de datos del mundo real, especialmente las series temporales, frecuentemente contienen valores faltantes (\textit{missing data}). Abordar incorrectamente estos datos puede introducir sesgos significativos en los análisis o reducir drásticamente la cantidad de datos utilizables para la modelización \cite[Capítulo 8]{Bishop2006}. La \textbf{imputación de datos} es el proceso de estimar y rellenar estos valores ausentes, con el objetivo de restaurar la completitud del conjunto de datos y mitigar los efectos adversos de la información faltante. La elección y aplicación de una técnica de imputación adecuada es, por tanto, un paso clave para la validez de cualquier análisis o modelo posterior construido sobre los datos.

% --- SECCIÓN 2: DEFINICIÓN FORMAL DEL PROBLEMA Y OBJETIVOS ---
\section{Definición del Problema de Imputación Experimental y Evaluación}
\label{sec:impl_problem_formal_revised}

El problema abordado en este trabajo consiste en desarrollar y evaluar funciones de imputación $f$ que, aplicadas al conjunto de datos de prueba con los valores faltantes, denotado como $(\mathcal{X}_{\text{missing}}^{\text{test}}, \mathbf{M}^{\text{test}})$, generen una estimación completada del conjunto de datos, $\hat{\mathcal{X}}^{\text{test}}$:
\[ \hat{\mathcal{X}}^{\text{test}} = f(\mathcal{X}_{\text{missing}}^{\text{test}}, \mathbf{M}^{\text{test}}) \]
Para los métodos de imputación basados en modelos predictivos, $f_{\theta}$, la función de imputación $f$ depende de un conjunto de parámetros $\theta^*$ que son optimizados durante una fase de entrenamiento previa.

El \textbf{objetivo principal de la evaluación} es cuantificar la precisión con la que las estimaciones $\hat{\mathcal{X}}^{\text{test}}$ reconstruyen los valores verdaderos originales, $\mathcal{X}_{\text{orig}}^{\text{test}}$. Para ello, se elimina deliberadamente un porcentaje de los valores observados del conjunto de datos, estableciendo así una categorización entre \textbf{datos faltantes naturales} (aquellos que están ausentes de forma inherente en el conjunto de datos original) y \textbf{datos faltantes artificiales} (aquellos que se eliminan intencionalmente para propósitos de evaluación).
La evaluación se centra exclusivamente en aquellas posiciones donde los datos faltantes fueron eliminados artificialmente, según lo indicado por la máscara $\mathbf{M}_{\text{artificial}}^{\text{test}}$, permitiendo una comparación directa entre los valores verdaderos conocidos y las estimaciones generadas por el método de imputación. 

\[
\text{Error} = \text{Metric}(\hat{\mathcal{X}}^{\text{test}}, \mathcal{X}_{\text{orig}}^{\text{test}}, \mathbf{M}_{\text{artificial}}^{\text{test}})
\]

Se busca identificar la función de imputación $f$ que minimice este error de reconstrucción, medido a través de diversas métricas como el Error Absoluto Medio (MAE), el Error Cuadrático Medio (MSE), la Raíz del Error Cuadrático Medio (RMSE) y el Error Relativo Medio (MRE), en el conjunto de datos de prueba.

% --- SECCIÓN 3: TAXONOMÍA GENERAL DE MÉTODOS DE IMPUTACIÓN 

\section{Métodos de Imputación en Series Temporales}
\label{sec:taxonomia_imputacion}

Las estrategias de imputación para series temporales varían ampliamente según la información que utilizan y los supuestos subyacentes sobre la naturaleza de los datos faltantes y la propia serie temporal \citep{Little2019statistical, Schafer1997analysis}. En este estudio, adoptamos una taxonomía que clasifica los métodos según el alcance de la información utilizada y la complejidad del modelo subyacente. Sea una serie temporal univariante $X = (x_1, x_2, \ldots, x_T)$, donde algunos $x_t$ pueden estar faltantes. Denotaremos el valor imputado para $x_t$ como $\hat{x}_t$.

\begin{itemize}
    \item \textbf{Métodos de Imputación con Información Local (Basados en Propagación o Interpolación)}:
    Estos métodos estiman los valores faltantes $x_t$ basándose predominantemente en los puntos de datos temporalmente adyacentes. Explotan la autocorrelación temporal, partiendo del supuesto de que las observaciones cercanas en el tiempo tienden a presentar valores similares \citep{Molenberghs2015handbook}. Son computacionalmente eficientes pero pueden ser imprecisos si los huecos son grandes o la serie es muy volátil.

        \begin{itemize}
            \item \textit{\glsp{ffill}} o \glsp{locf}: Si $x_t$ es faltante, se imputa con el último valor observado que lo precede. Sea $t_0 < t$ el índice temporal más reciente tal que $x_{t_0}$ está observado. Entonces:
            \begin{equation}
            \hat{x}_t = 
            \begin{cases}
            x_{t_0} & \text{si } \{\tau \mid \tau < t \text{ y } x_\tau \ne \text{NaN}\} \text{ no es vacío} \\
            0 & \text{en caso contrario}
            \end{cases}\label{eq:locf}
            \end{equation}

            donde $t_0 = \max\{\tau \in \{1, \dots, T\} \mid \tau < t \text{ y } x_\tau \ne \text{NaN}\}$.
            
            Este método es simple y comúnmente usado, especialmente en datos médicos o financieros donde se asume persistencia del estado \citep{Shao1992handling}.

            \item \textit{\glsp{bfill}} o \glsp{nocb}: Si $x_t$ es faltante, se imputa utilizando el siguiente valor conocido en la secuencia. Sea $t_1 > t$ el índice temporal más cercano tal que $x_{t_1}$ está observado. Entonces:
            \begin{equation}
            \hat{x}_t = 
            \begin{cases}
            x_{t_1} & \text{si } \{\tau \mid \tau > t \text{ y } x_\tau \ne \text{NaN}\} \text{ no es vacío} \\
            0 & \text{en caso contrario}
            \end{cases}\label{eq:nocb}
            \end{equation}

            donde $t_1 = \min\{\tau \in \{1, \dots, T\} \mid \tau > t \text{ y } x_\tau \ne \text{NaN}\}$.
            
            Utilizado cuando se espera que el futuro inmediato influya en el valor actual \citep{Allison2001missing}.

            \item \textit{Interpolación Lineal}: Si $x_t$ es faltante y se encuentra entre dos valores observados $x_{t_a}$ (anterior, $t_a < t$) y $x_{t_p}$ (posterior, $t_p > t$), con todos los valores entre $x_{t_a}$ y $x_{t_p}$ (excluyendo $x_t$) también faltantes o siendo estos los más cercanos observados, se asume una trayectoria lineal. El valor imputado $\hat{x}_t$ se calcula como:
            \begin{equation}
            \hat{x}_t = x_{t_a} + (x_{t_p} - x_{t_a}) \frac{t - t_a}{t_p - t_a}
            \label{eq:linear_interpolation}
            \end{equation}
            Este método es una mejora sobre LOCF/NOCB si hay fluctuaciones entre los puntos conocidos \citep{Hastie2009}. Para huecos al inicio o final de la serie, se puede recurrir a extrapolación o a LOCF/NOCB.
        \end{itemize}
        Estos métodos son generalmente eficientes para rellenar huecos cortos, pero su rendimiento puede decaer ante lagunas de datos más extensas o patrones no lineales.

    \item \textbf{Métodos de Imputación Estacionarios (Basados en Estadísticas Globales o de Ventana)}:
    Se fundamentan en el supuesto de que ciertas propiedades estadísticas de la serie temporal (o segmentos de ella), como la media o la mediana, son representativas de los valores faltantes. La \textbf{estacionariedad} implica que las propiedades estadísticas de la serie (media, varianza, autocorrelación) son invariantes en el tiempo \citep{Box2015time}. Estos métodos utilizan estadísticas calculadas sobre los datos observados $X_{obs}$.

        \begin{itemize}
            \item \textit{Imputación por Media}: Reemplaza los valores faltantes $x_t$ con la media aritmética de todos los valores observados en la serie (o de una ventana temporal relevante $W$ alrededor de $t$ si se aplica localmente):
            \begin{equation}
            \hat{x}_t = \frac{1}{|X_{\text{missing}}^{test}|} \sum_{x_i \in X_{\text{missing}}^{test} } x_i
            \end{equation}
            Este método es simple pero puede reducir la varianza de la serie e ignorar la estructura temporal y las relaciones con otras variables \citep{Donders2006review}.

            \item \textit{Imputación por Mediana}: Similar al anterior, pero utiliza la mediana de los valores observados, lo que lo hace más robusto a valores atípicos (outliers) \citep{Enders2010applied}:
            \begin{equation}
            \hat{x}_t = \text{mediana}(X_{\text{missing}})
            \label{eq:median_imputation}
            \end{equation}
            Comparte limitaciones similares a la imputación por media en cuanto a la distorsión de la estructura de la serie.
        \end{itemize}
        Aunque simples, estos métodos tienden a subestimar la varianza y aplanar la estructura temporal de la serie, y son generalmente inapropiados si la serie no es estacionaria o si los datos faltantes no son completamente al azar.

    \item \textbf{Métodos de Imputación Predictivos (Basados en Modelos)}:
    Emplean modelos estadísticos o de \glsentryshort{ml} para predecir los valores faltantes, aprendiendo patrones y relaciones a partir de los datos observados. Son capaces de capturar dependencias más complejas \citep{AbuMostafa2012}. El objetivo es aprender una función $f$ tal que $\hat{x}_t = f(X_{missing}^{test}, M^{test})$

        \begin{itemize}
            \item \textit{Modelos de Regresión}:
            En este contexto, la función $f$ representa el modelo de regresión (lineal, no lineal, etc.) ajustado sobre los datos observados. Para imputar un valor faltante $\hat{x}_t$, $f$ toma como entradas los valores observados de otras variables.
            La función $f$ se traduce en la ecuación del modelo de regresión. Por ejemplo, para una regresión lineal multivariante, donde $\hat{x}_t$ es la variable dependiente y $Z_t = (z_{t,1}, \ldots, z_{t,p})$ son las variables predictoras observadas en el instante $t$ (extraídas de $\mathcal{X}_{missing}^{test}$):
            
            \begin{equation}
            \hat{x}_t = \beta_0 + \sum_{j=1}^p \beta_j z_{t,j}
            \end{equation} 
            
            Los coeficientes $\beta_j$ son parámetros aprendidos por el modelo de regresión. La máscara $\mathbf{M}^{test}$ es fundamental para identificar los $x_t$ que deben ser imputados y para asegurar que solo se utilicen las variables predictoras disponibles y válidas en $\mathcal{X}_{missing}^{test}$.

            \item \textit{Modelos Autorregresivos (e.g., ARIMA)}:
            En los modelos autorregresivos, la función $f$ captura la dependencia de un valor $x_t$ con sus propios valores pasados. La función $f$ representa el modelo autorregresivo (AR, MA, ARMA, ARIMA) con los parámetros estimados a partir de la serie observada.
            Para imputar $\hat{x}_t$, $f$ utiliza los valores históricos observados de la misma serie. Por ejemplo, para un modelo autorregresivo AR(p), la función $f$ se materializa como:
            
            \begin{equation}
            \hat{x}_t = c + \sum_{i=1}^p \phi_i x_{t-i} + \epsilon_t
            \end{equation}

            donde los $x_{t-i}$ son los valores de la serie en instantes anteriores (observados y extraídos de $\mathcal{X}_{missing}^{test}$), los $\phi_i$ son los coeficientes autorregresivos aprendidos, $c$ es una constante y $\epsilon_t$ es un término de error. La máscara $\mathbf{M}^{test}$ informa a $f$ qué valores de $x$ están realmente observados y cuáles necesitan imputación, y el modelo manejaría internamente la situación si alguno de los $x_{t-i}$ también fuera faltante (por ejemplo, mediante un proceso iterativo o la imputación previa de esas dependencias).

           \item \textit{Modelos Basados en Matrices}:
            Estos modelos transforman la serie temporal en una matriz (e.g., Hankel, Toeplitz) y luego aplican técnicas de completado de matrices para imputar los valores faltantes.
    
            En el caso de la Imputación Hankel (HI) \citep{PoudevigneJones24Hankel} , la serie temporal $x(1), \ldots, x(n)$ en $\mathbb{R}^d$ se reestructura en una matriz de Hankel por bloques, $H$, de tamaño $(n-k+1) \times (kd)$, donde $k$ es el retardo (lag) fijo y $d$ la dimensión de la serie (para series univariadas, $d=1$). Esta matriz $H$ es construida de tal manera que cada fila representa un segmento de la serie temporal:
    
            \[
            H = \begin{pmatrix}
            x(1) & x(2) & \cdots & x(k) \\
            x(2) & x(3) & \cdots & x(k+1) \\
            \vdots & \vdots & \ddots & \vdots \\
            x(n-k+1) & x(n-k+2) & \cdots & x(n)
            \end{pmatrix}
            \]
    
            La tarea de imputación se reformula como un problema de completado de matriz de bajo rango. El objetivo es encontrar una matriz $X$ de bajo rango que coincida con los valores observados de $H$ en las posiciones conocidas. El problema se formula como una minimización de la norma nuclear de la matriz $X$ (sumatoria de sus valores singulares, $\|X\|_*$), sujeta a las restricciones de que los valores observados en las posiciones $\Omega$ deben ser preservados. La norma nuclear es una relajación convexa de la función de rango:
    
            \begin{equation}
            \min_X \|X\|_* \quad \text{s.t.} \quad X(i,j) = M(i,j) \quad \forall (i,j) \in \Omega
            \end{equation}
            Donde $M$ es la matriz Hankel original con valores faltantes, y $\Omega$ es el conjunto de índices de los valores observados (es decir, las entradas conocidas de $M$). La matriz $X$ obtenida al resolver este problema es la matriz Hankel completada.
    
            Los valores imputados $\hat{x}_t$ corresponden a la reconstrucción de la serie temporal a partir de esta matriz $X$ completada. El valor imputado $\hat{x}_t$ se obtiene usualmente promediando las entradas correspondientes en la matriz $X$ completada.
    
            La máscara $\mathbf{M}^{test}$ en este contexto se traduce en el conjunto $\Omega$, que identifica las entradas observadas de la matriz Hankel que deben ser preservadas durante el proceso de completado. Las entradas fuera de $\Omega$ son las que el algoritmo de completado de matriz calculará para imputar los valores faltantes en la serie temporal subyacente.
    
            
            \item \textit{Modelos de Aprendizaje Profundo}: Arquitecturas avanzadas como las Redes Neuronales Recurrentes (RNNs, e.g., LSTM \citep{Hochreiter1997lstm}, \glsp{gru} \citep{Cho2014gru}) o los Transformers (e.g., SAITS \citep{Du_saits_2023}) son especialmente adecuadas para series temporales. Dada su capacidad para modelar dependencias temporales complejas y de largo alcance, así como interacciones no lineales entre múltiples series, pueden aprender representaciones ricas de los datos observados para inferir los valores faltantes \citep[Capítulos 10 y 12]{Goodfellow2016}. Por ejemplo, un modelo podría ser entrenado para predecir $x_t$ dado un contexto de valores observados $X_{missing}^{test} = (x_{t-k}, \ldots, x_{t-1}, x_{t+1}, \ldots, x_{t+k})$.
            
        \end{itemize}
        Estos métodos son generalmente más potentes y pueden proporcionar imputaciones más precisas, especialmente en presencia de patrones complejos o cuando el mecanismo de generación de los datos faltantes depende de otras variables observadas o no observadas. 
        \citep{Rubin1976inference, Little2019statistical}. 
        \begin{equation}
        f_{\theta}(X^{\text{test}}_{\text{missing}}, M^{\text{test}}_{\text{missing}}) = \text{Transformer}_\theta({X_{missing}^{test}}, M_{missing}^{test})
        \end{equation}
        Específicamente, pueden ser más efectivos en escenarios donde los datos son \textbf{Faltantes Aleatoriamente (MAR, Missing At Random)} –es decir, la probabilidad de que un valor falte depende solo de otras variables observadas– o incluso \textbf{Faltantes No Aleatoriamente (MNAR, Missing Not At Random)} –donde la probabilidad de ausencia depende del propio valor que falta o de otras variables no observadas–, siempre y cuando el modelo sea capaz de capturar o se le proporcione información sobre dicho mecanismo de ausencia 
        Sin embargo, estos enfoques predictivos pueden ser computacionalmente más costosos, requerir conjuntos de datos más grandes para un entrenamiento efectivo y son susceptibles al sobreajuste si no se regularizan adecuadamente.
\end{itemize}
La elección del método de imputación más apropiado depende de múltiples factores, incluyendo las características de los datos (univariante/multivariante, longitud, estacionalidad, tendencia), el mecanismo, patrón y proporción de valores faltantes, y los recursos computacionales disponibles \citep{Rubin1976inference, VanBuuren2018flexible}.


% --- SECCIÓN 5: MÉTRICAS DE EVALUACIÓN ---
\section{Métricas de Evaluación para los Métodos de Imputación}
\label{sec:evaluation_metrics_hadamard}

Para cuantificar el rendimiento de los métodos de imputación, se emplean métricas de error estándar como el Error Cuadrático Medio (MSE), la Raíz del Error Cuadrático Medio (RMSE), el Error Absoluto Medio (MAE) y el Error Relativo Medio (MRE). Estas métricas, ya introducidas en la \Cref{subsec:Metricas}, requieren una adaptación específica para el problema de datos faltantes. La evaluación del error se realiza \textit{exclusivamente} en las posiciones donde los datos originales fueron artificialmente eliminados. Para ello, comparamos los valores imputados $\hat{\mathcal{X}}^{\text{test}}$ con los valores verdaderos $\mathcal{X}_{\text{orig}}^{\text{test}}$ únicamente en las posiciones indicadas por la máscara $\mathbf{M}_{\text{artificial}}^{\text{test}}$. Estas adaptaciones se formalizan matemáticamente a continuación.

Sean:
\begin{itemize}
    \item $\hat{\mathcal{X}} = \hat{\mathcal{X}}^{\text{test}} \in \mathbb{R}^{N \times T \times D}$: Tensor de valores imputados (conjunto de prueba).
    \item $\mathbf{X}_{\text{orig}} = \mathcal{X}_{\text{orig}}^{\text{test}} \in \mathbb{R}^{N \times T \times D}$: Tensor de valores verdaderos originales (conjunto de prueba).
    \item $\mathbf{M}_{\text{ind}} = \mathbf{M}_{\text{artificial}}^{\text{test}} \in \{0, 1\}^{N \times T \times D}$: Máscara binaria que indica las posiciones donde los valores fueron artificialmente eliminados y posteriormente imputados ($1$ si fue imputada, $0$ si no).
\end{itemize}

Las métricas se definen como:

\begin{definition}[Error Absoluto Medio (MAE)]
\label{def:mae_hadamard}
El Error Absoluto Medio (MAE) se calcula como:
\begin{equation}
\text{MAE}(\hat{\mathcal{X}}, \mathbf{X}_{\text{orig}}, \mathbf{M}_{\text{ind}}) = \frac{||\lvert \hat{\mathcal{X}} - \mathbf{X}_{\text{orig}} \rvert \odot \mathbf{M}_{\text{ind}}||_1}{||\mathbf{M}_{\text{ind}}||_1}
\label{eq:mae_hadamard}
\end{equation}
según \citep{Willmott2005} donde $\lvert \cdot \rvert$ denota el valor absoluto elemento a elemento, $\odot$ es el producto de Hadamard (Ver \Cref{def:hadamard_product_general}), y $||\cdot||_1$ es la norma $\ell_1$ (suma de los valores absolutos de los elementos). El término $\odot \mathbf{M}_{\text{ind}}$ asegura que solo se consideren los errores en las posiciones imputadas. El denominador normaliza por el número total de puntos imputados.
\end{definition}


\begin{definition}[Raíz del Error Cuadrático Medio (RMSE)]
\label{def:rmse_hadamard}
La Raíz del Error Cuadrático Medio (RMSE) se define por:
\begin{equation}
\text{RMSE}(\hat{\mathcal{X}}, \mathbf{X}_{\text{orig}}, \mathbf{M}_{\text{ind}}) = \sqrt{\frac{||(\hat{\mathcal{X}} - \mathbf{X}_{\text{orig}})^{\odot 2} \odot \mathbf{M}_{\text{ind}}||_1}{||\mathbf{M}_{\text{ind}}||_1}}
\label{eq:rmse_hadamard}
\end{equation}

según \citep{Chai2014}. 
Aquí, $(\hat{\mathcal{X}} - \mathbf{X}_{\text{orig}})^{\odot 2}$ representa el error cuadrático elemento a elemento. Similar al MAE, la máscara $\mathbf{M}_{\text{ind}}$ selecciona los errores cuadráticos en las posiciones relevantes, se suman, se promedian y finalmente se toma la raíz cuadrada.
\end{definition}

\begin{definition}[Error Relativo Medio (MRE)]
\label{def:mre_hadamard}
El Error Relativo Medio (MRE) se calcula como:
\begin{equation}
\text{MRE}(\hat{\mathcal{X}}, \mathbf{X}_{\text{orig}}, \mathbf{M}_{\text{ind}}) = \frac{||\lvert \hat{\mathcal{X}} - \mathbf{X}_{\text{orig}} \rvert \odot \mathbf{M}_{\text{ind}}||_1}{||\lvert \mathbf{X}_{\text{orig}} \rvert \odot \mathbf{M}_{\text{ind}}||_1 + \epsilon}
\label{eq:mre_hadamard}
\end{equation}

según \citep{Chai2014rmse}.
El numerador corresponde a la suma de los errores absolutos enmascarados. El denominador es la suma de los valores absolutos verdaderos originales en esas mismas posiciones enmascaradas. Se añade una pequeña constante $\epsilon > 0$ al denominador para prevenir la división por cero en caso de que algún valor original en una posición imputada sea cero.
\end{definition}
Adicionalmente, se reporta el Error Cuadrático Medio (MSE), que es simplemente el cuadrado del RMSE antes de tomar la raíz cuadrada.

%%%%SAITS

\section{Aprendizaje Profundo para la Imputación en Series Temporales}
\label{subsubsec:ventajas_saits_sintetizado}

Las redes neuronales profundas son cada vez más usadas para la imputación en series temporales, superando a métodos tradicionales. Enfoques basados en \glsentryshort{rnn} como BRITS \citep{Cao2018brits}, GRU-D \citep{Che2018recurrent}, o \glsentryshort{gan} para series temporales \citep{Luo2018multivariate_e2gan} han mostrado resultados prometedores. Sin embargo, arquitecturas basadas en autoatención, como el Transformer \citep{Vaswani2017} y modelos derivados como SAITS \citep{Du2023SAITS} para imputación, presentan ventajas distintivas:

\begin{itemize}
    \item \textbf{Modelado Eficaz de Dependencias Temporales a Largo Plazo:}
    Los modelos basados en autoatención como SAITS destacan por su capacidad para modelar dependencias a largo plazo, superando una limitación clave de las \glsentryshort{rnn} (e.g., LSTM \citep{Hochreiter1997lstm}, \glsentryshort{gru} \citep{Cho2014gru}). El procesamiento secuencial de las \glsentryshort{rnn} puede llevar a la pérdida de información o al desvanecimiento de gradientes en secuencias largas, obstaculizando la captura de relaciones distantes \citep{Bengio1994learning, Pascanu2013difficulty}. En cambio, la autoatención ofrece un acceso directo entre todos los puntos de la secuencia (camino de $O(1)$ vs. $O(L)$ en RNNs), lo que es fundamental para imputar valores que dependen de información lejana \citep{Vaswani2017}.
    \item \textbf{Captura Integral de Dependencias Conjuntas (Inter e Intra-Variable):}
    En series multivariantes, los modelos de autoatención pueden modelar holísticamente tanto la evolución individual de cada variable (intra-variable) como las interacciones entre diferentes variables (inter-variable) a lo largo del tiempo. Al permitir que cada variable en cada instante atienda a todas las demás en todos los instantes, se pueden aprender patrones de correlación complejos y dinámicos que modelos más simples podrían omitir.

    \item \textbf{Flexibilidad ante Patrones Complejos sin Supuestos Restrictivos:}
    Muchos métodos tradicionales requieren supuestos como la \textbf{estacionariedad} (propiedades estadísticas constantes en el tiempo \citep{Box2015time}). Los modelos basados en Transformers, como arquitecturas de \glsentryshort{dl}, aprenden patrones no lineales y no estacionarios directamente de los datos, adaptándose a dinámicas temporales diversas sin necesidad de preprocesamientos extensivos o supuestos restrictivos sobre la estructura de los datos.

    \item \textbf{Potencial de Robustez frente a Ruido y Outliers:}
    El mecanismo de atención, al ponderar la importancia de diferentes segmentos de la secuencia, tiene el potencial de asignar pesos menores a partes ruidosas o con outliers. Esto podría llevar a imputaciones más robustas, aunque la efectividad de esta propiedad depende del aprendizaje adecuado de dichas ponderaciones.
\end{itemize}

Frente a los modelos mencionados inicialmente y considerando el conjunto de ventajas expuestas, para este trabajo hemos seleccionado SAITS \citep{Du2023SAITS}. La principal razón de esta elección radica en su demostrada capacidad para el \textbf{modelado eficaz de dependencias temporales a largo plazo} (como se detalló en el primer punto), una característica esencial para la imputación precisa en series temporales complejas donde la información relevante puede estar significativamente separada en el tiempo. Esta fortaleza, combinada con su habilidad para manejar dependencias multivariantes y su flexibilidad, lo posiciona como un candidato ideal para los desafíos de imputación abordados en este estudio.

% --- FIN DEL CAPÍTULO DE METODOLOGÍA ---